% Introduction, drafted 2026-09-05; prose revised 2026-09-05 for the SaTML submission.
\section{Introduction}\label{sec:intro}

Language models reproduce text they were trained on, and some of that text is copyrighted. Near-access-freeness (NAF)~\cite{vyas2023provable} offers a mechanism-level remedy: a model whose output distribution stays within a divergence budget of a \emph{safe} model that never saw a protected work cannot, the argument goes, be much more likely than the safe model to reproduce it. Anchored Decoding~\cite{he2026anchored} is the first implementation of this idea at inference time for an arbitrary risky model. It fuses the risky model with a permissively trained anchor token by token, charges each step's KL divergence against a banked allowance, and certifies a sequence-level budget of $K = k\,T_{\max}$ nats per prompt. The certificate holds for any risky model, is enforced during decoding, and comes with a proof.

This paper asks what the certificate is worth to the party it is meant to protect. A rights-holder cares about an event, the reproduction of a specific passage, and about a system that answers many prompts. The certificate bounds an expectation, for one prompt, in KL divergence. We audit the distance between the two in the way the differential-privacy community audits a claimed $\varepsilon$~\cite{jagielski2020auditing,nasr2023tight,steinke2023privacy}: by constructing measurements that could contradict it. Every audit step below has a failure mode, and we report which ones fired.

We first establish that the accounting is right. On all 26{,}999 trajectories of the released Anchored Decoding runs~\cite{vijayavallabh2026audit}, recomputed in single precision, the per-step and per-trajectory invariants hold with a largest overshoot of zero (Section~\ref{sec:accounting}). Everything reported afterwards therefore concerns the guarantee rather than a defect in the implementation. We then measure what the guarantee means at the budgets of the released runs. At $k=3$ and $k=5$ the constraint is active on $0.2\%$ and $0.1\%$ of decode steps, the risky model is served unchanged on $99\%$, and $57\%$ of generations are identical between the two budgets, so the mechanism is close to the identity there (Section~\ref{sec:regime}). For the 758 CopyBench passages under the mechanism's anchor, the bound that the certificate implies on the probability of reproducing a passage is vacuous for every passage at $k \ge 3$ and for $44\%$ of passages at $k=1$; it becomes informative only at $k \le 0.5$, where the anchor also writes the first tokens of every answer (Section~\ref{sec:certificate}). Because the certificate bounds an expectation, individual outputs can exceed it: at $k=1$ the realised log-likelihood ratio of an output exceeds $K$ for $8\%$ of trajectories across six prompt classes, and single tokens are served with $e^{26}$ times their anchor probability (Section~\ref{sec:tails}).

The remaining question is whether any of this is exploitable. The audited risky model, Llama-3.1-8B-Instruct, has not memorised the passages, so we build one that has by fine-tuning it on the CopyBench passages until greedy decoding reproduces them, and we attack the deployed decoder with it under a black-box threat model (Section~\ref{sec:attacks}). The attacks separate two regimes. At $k \le 3$ the budget binds against the memorising model (utilisation above $0.8$) and a single query recovers about $1\%$ of a passage, even though the certificate for that passage is vacuous; the protection is real, but the certificate does not describe it. Between $k=5$ and $k=20$, budgets that the mechanism's authors also evaluate, the protection disappears. A single query recovers $48\%$ of a passage at $k=20$, nearly as much as the unconstrained model does, and windowed queries in the spirit of Cohen's composition theorem~\cite{cohen2025blameless}, each within its own budget, recover $86\%$. At the intermediate budgets composition doubles what a single query extracts. The prefix debt that the mechanism charges on a memorised prompt turns out to be its most effective component against short windows, which the original analysis does not exploit.

\paragraph{Contributions}
(C1)~An exact accounting audit of a deployed NAF mechanism on $5.2$ million logged decode steps, with zero violations. (C2)~A regime audit showing that the budgets of the released runs, $k \in \{3, 5\}$, leave the mechanism inactive on $99\%$ of steps, together with a small-budget sweep against the risky model alone and the anchor alone. (C3)~A certificate-strength audit that converts the KL budget into a per-passage cap on the reproduction probability and shows the cap is vacuous for all CopyBench passages at $k \ge 3$. (C4)~An event-level audit of the realised log-likelihood ratio, the privacy-loss analogue of the certificate, with anytime-valid confidence sequences replacing fixed-sample bounds. (C5)~Adaptive attacks with a memorising risky model, including the first empirical instantiation of the composition theorem against an inference-time NAF mechanism, with every query within budget. (C6)~An audit protocol that can fail, released with code, logs, and prompt sets so that deployers and rights-holders can run it against other anchors, budgets, and passages.
